\begin{thema}{Γ}
Στο σχήμα δίνονται οι γραφικές παραστάσεις $C_1$ και $C_2$ δύο συναρτήσεων, μιας συνεχούς συνάρτησης $f:[-1, 4]\to\mathbb{R}$ και της παραγώγου της $f'$. Γνωρίζουμε ότι η $C_1$ διέρχεται από τα σημεία $(-1, 0)$, $(1, 3)$ και $(4, -1)$, ενώ η $C_2$ διέρχεται από τα σημεία $(-1, 3)$, $(1, 0)$ και $(4, -2)$.

\begin{center}
\begin{tikzpicture}
\begin{axis}[
    width=11cm, height=8.5cm,
    axis lines=middle,
    xmin=-1.5, xmax=4.5,
    ymin=-2.5, ymax=3.5,
    xtick={-1,1,2,3,4},
    ytick={-2,-1,1,2,3},
    xlabel={$x$},
    ylabel={$y$},
    grid=both,
    grid style={line width=.1pt, draw=gray!20},
    major grid style={line width=.2pt, draw=gray!40},
    enlargelimits=false,
    ticklabel style={font=\footnotesize},
    samples=100
]

% C1 (f) on [-1, 1]
\addplot[very thick, maincolor, domain=-1:1] {-0.75*x^2 + 1.5*x + 2.25};
% C1 (f) on [1, 4]
\addplot[very thick, maincolor, domain=1:4] {(2/27)*(x-1)^3 - (2/3)*(x-1)^2 + 3} node[pos=0.8, above right] {$C_1$};

% C2 (f') on [-1, 1]
\addplot[very thick, secondcolor, domain=-1:1] {-1.5*x + 1.5};
% C2 (f') on [1, 4]
\addplot[very thick, secondcolor, domain=1:4] {(2/9)*(x-1)^2 - (4/3)*(x-1)} node[pos=0.8, below left] {$C_2$};

% Points for C1
\addplot[mark=*, maincolor, only marks, mark size=1.5pt] coordinates {(-1,0) (1,3) (4,-1)};
\node[maincolor, above, font=\footnotesize] at (axis cs:1,3) {$(1,3)$};
\node[maincolor, below left, font=\footnotesize] at (axis cs:4,-1) {$(4,-1)$};

% Points for C2
\addplot[mark=*, secondcolor, only marks, mark size=1.5pt] coordinates {(-1,3) (1,0) (4,-2)};
\node[secondcolor, right, font=\footnotesize] at (axis cs:-1,3) {$(-1,3)$};
\node[secondcolor, below right, font=\footnotesize] at (axis cs:4,-2) {$(4,-2)$};

\end{axis}
\end{tikzpicture}
\end{center}

\begin{erwthma}
  \item Να αιτιολογήσετε ποια από τις καμπύλες $C_1, C_2$ παριστάνει τη συνάρτηση $f$ και ποια την παράγωγό της $f'$. Στη συνέχεια να βρείτε τα διαστήματα μονοτονίας και τα ακρότατα της $f$.
  \item Να αποδείξετε ότι η $f$ είναι κοίλη στο $[-1, 4]$ και ότι η εξίσωση $f(x) = 0$ έχει ακριβώς μία ρίζα στο διάστημα $(1, 4)$.
  \item Να αποδείξετε ότι $\dint_{-1}^1 f(x)\,\d x > 3$ και να υπολογίσετε το εμβαδόν του χωρίου που περικλείεται από την $C_2$, τον άξονα $x'x$ και τις ευθείες $x=-1$ και $x=4$.
  \item Να υπολογίσετε το ολοκλήρωμα $\dint_1^4 x f''(x)\,\d x$.
\end{erwthma}
\end{thema}
